%==============================================================================
% ICSE 2027 -- Tool Demonstration and Data Showcase Track
% arcade-agent: Architecture Recovery and Conformance Inside the Agent's Loop
%
% SKELETON ONLY. Section scaffolding, verified tables, and figure stubs.
% The prose is not written; each section carries a \secbudget and TODO notes
% taken from OUTLINE.md in this directory.
%
% BEFORE DRAFTING, verify against the official CFP
%   https://conf.researchr.org/track/icse-2027/icse-2027-demonstrations
% the four constraints this file assumes:
%   (1) 4 pages INCLUSIVE of references, figures, tables, appendices
%   (2) IEEEtran, [10pt,conference], no compsoc
%   (3) single-anonymous -- authors ARE named
%   (4) the demo video URL is appended to the END of the abstract
% All four came from search snippets, not the CFP itself. See OUTLINE.md Sec. 1.
%==============================================================================

\documentclass[10pt,conference]{IEEEtran}
% Do NOT add compsoc or compsocconf.

\usepackage{cite}
\usepackage{graphicx}
\usepackage{booktabs}
\usepackage{xcolor}
\usepackage{listings}
\usepackage{amssymb}
\usepackage{balance}          % balance the last page's columns
\usepackage[hidelinks]{hyperref}
\usepackage{url}
\usepackage{tikz}
\usetikzlibrary{arrows.meta,positioning,fit,backgrounds}

%------------------------------------------------------------------------------
% Draft scaffolding. Flip \draftmodetrue -> \draftmodefalse for the submission
% PDF: every \todo and \secbudget disappears and the page count becomes real.
%------------------------------------------------------------------------------
\newif\ifdraftmode
\draftmodetrue

\ifdraftmode
  \newcommand{\todo}[1]{{\color{red}\small[\textbf{TODO} #1]}}
  \newcommand{\secbudget}[1]{{\color{blue!70!black}\scriptsize\textit{[budget: #1]}}\par}
\else
  \newcommand{\todo}[1]{}
  \newcommand{\secbudget}[1]{}
\fi

% Consistent typesetting for tool names, so a global rename is one edit.
\newcommand{\tool}[1]{\texttt{#1}}
\newcommand{\arcade}{\textsc{arcade-agent}}
\newcommand{\arcguard}{\textsc{arcade-guard}}

\lstset{
  basicstyle=\ttfamily\scriptsize,
  breaklines=true,
  columns=fullflexible,
  frame=single,
  framesep=3pt,
  aboveskip=4pt,
  belowskip=4pt,
  xleftmargin=2pt,
  literate={->}{{$\rightarrow$}}1
}

\begin{document}

%==============================================================================
% TITLE AND AUTHORS
% Single-anonymous review: real names and affiliations go here.
%==============================================================================
\title{\arcade{}: Architecture Recovery and Conformance\\Inside the Coding Agent's Loop}

\author{
  \IEEEauthorblockN{\todo{author names -- confirm list and order}}
  \IEEEauthorblockA{\todo{affiliations}\\\todo{emails}}
}

\maketitle

%==============================================================================
% ABSTRACT -- target <= 150 words.
% The video URL is a submission requirement and goes at the very END.
%==============================================================================
\begin{abstract}
\todo{$\le$150 words, one sentence per slot:}
\todo{(1) Coding agents produce working code and decaying architecture; nothing
guards architecture while code is being written.}
\todo{(2) \arcade{} is an installable library and MCP server that puts architecture
recovery, smell detection, decay metrics, and conformance checking inside the
agent's loop.}
\todo{(3) Mechanism: 19 MCP tools over 7 language parsers and 5 recovery
algorithms, with per-call token budgeting and session handles so an agent can
afford to call them.}
\todo{(4) The guardrail is proactive -- it answers where code should live and
what it may depend on before the code exists.}
\todo{(5) Evidence: on greenfield tasks an unaided agent introduced a forbidden
dependency in 3 of 14 runs; with the guardrail, 0 of 14, with no false positives.}
\todo{(6) Availability: pip, MIT, public repo.}

\medskip
\noindent Demonstration video: \url{TODO-video-url}
\todo{REQUIRED by the track; must be live and reachable without login before
submission. See OUTLINE.md Sec. 6.}
\end{abstract}

\begin{IEEEkeywords}
software architecture, architecture recovery, architectural decay, conformance
checking, AI coding agents, Model Context Protocol
\end{IEEEkeywords}

%==============================================================================
\section{Introduction}
\secbudget{0.8 col}
%==============================================================================
\todo{Hook. Type checkers guard types, linters guard style, tests guard
behavior -- nothing guards architecture while the code is being written. The
framing already exists in arcade-analyze-skill/GUARDRAIL\_PLAN.md; tighten it,
do not re-invent it.}

\todo{Why it matters: architectural decay is measured and costly. Cite the decay
and change-proneness line~\cite{le2018decay,le2021decay-predictor,le2015change}.}

\todo{Why existing tools do not fit the agent loop: batch, slow, human-facing,
output measured in megabytes of HTML. An agent has a token budget and a latency
budget. Post-hoc CI conformance tells the agent it was wrong after the fact, and
a red X with no fix makes an agent thrash.}

\todo{State the claim in one sentence -- see OUTLINE.md Sec. 2. Then be explicit
about what is NOT claimed: the recovery algorithms are prior
work~\cite{tzerpos2000acdc,garcia2011arc,andritsos2004limbo,maqbool2004wca,laser2020arcade}.
The contribution is the delivery vehicle. Saying this plainly is cheaper than
having a reviewer discover it.}

\noindent\textbf{Contributions.}
\begin{itemize}
  \item \todo{An installable, framework-agnostic architecture-analysis library
        (\tool{pip install arcade-agent}) with 7 language parsers and 5 recovery
        algorithms.}
  \item \todo{An MCP server exposing 19 tools with per-call \tool{max\_tokens}
        budgeting and a session-handle protocol, so agents consume summaries,
        not dumps.}
  \item \todo{Four task-shaped context tools that answer an agent's actual
        question instead of returning a graph.}
  \item \todo{An executable architecture contract plus proactive guardrail
        tools, with evidence that they eliminate a violation class an unaided
        agent produces roughly one run in five on greenfield code.}
\end{itemize}

%==============================================================================
\section{Motivating Scenario}
\label{sec:motivation}
\secbudget{0.8 col + Fig.~\ref{fig:scenario}}
%==============================================================================
\todo{One concrete vignette -- use the greenfield orders service, because it
doubles as the evaluated case in Sec.~\ref{sec:evidence}. An agent is told to
add order listing and totals in a repo whose API layer must not touch the store.
Unaided it imports app.store from app.api: the shortest path to a working
feature. With the guardrail it calls \tool{propose\_placement}, is told the logic
belongs in \tool{service}, calls \tool{preview\_impact} on the api->store edge,
is told it would violate, and routes through the service layer instead.}

\begin{figure}[t]
  \centering
  \fbox{\parbox[c][2.6cm][c]{0.92\columnwidth}{\centering
    \todo{FIGURE 1 -- two import graphs side by side: the unaided result with
    the forbidden \tool{api}$\rightarrow$\tool{store} edge in red, and the
    guarded result, clean. Draw in TikZ. Keep to one column. This figure is the
    whole paper in one picture.}}}
  \caption{The same task, with and without the guardrail.}
  \label{fig:scenario}
\end{figure}

%==============================================================================
\section{Architecture and Implementation}
\label{sec:impl}
%==============================================================================
\arcade{} is a library first and a server second. The analysis core is a
sequence of pure stages --- \tool{ingest}, \tool{parse}, \tool{recover},
\tool{detect\_smells}, \tool{compute\_metrics}, \tool{compare} --- each
independently callable, each returning a plain dataclass. The MCP server, the
CI entry points and the guardrail are all consumers of that core rather than
layers inside it (Fig.~\ref{fig:pipeline}).

\subsection{Parsing}
Each language is a \tool{LanguageParser} implementation registered through a
\tool{@register\_parser} decorator, so adding a language touches no shared
dispatch code. Seven are registered: Java and Python are core dependencies and
work out of the box; C, TypeScript, Go, Kotlin and Rust arrive with the
\tool{[languages]} extra. Parsing is two-pass --- a per-file AST walk that
depends only on that file's contents, then a link pass that resolves imports,
inheritance and qualified references across the whole file set.

Polyglot repositories are merged into one graph, but relinking is deliberately
\emph{family-scoped}. Java and Kotlin share the \tool{jvm} family and are
relinked against each other; every other language forms a single-member family
and is merged without cross-language edges. Unqualified name resolution is
likewise confined to a family, and only when the leaf name is unique within it,
so a Python \tool{Base} can never be mistaken for a Java one. Cross-family FQN
collisions are remapped rather than silently overwritten, and the count is
surfaced in the result's metadata. \todo{The JVM pair is the only combination
we have validated. Say so here rather than in a footnote --- an overclaim about
general cross-language recovery is the cheapest possible reject.}

\subsection{Recovery}
Five recovery algorithms are available: package-based grouping (\tool{pkg}, the
default), \tool{wca}, \tool{acdc}, and the LLM-assisted \tool{arc} and
\tool{limbo}. These are reimplementations of prior
work~\cite{tzerpos2000acdc,garcia2011arc,andritsos2004limbo,maqbool2004wca}, and
we claim no novelty for them.

The guardrail path, however, uses only \tool{pkg} and the spec's own glob
mapping. This is a deliberate restriction. A guardrail consulted dozens of times
in a single agent session must return the same verdict for the same code every
time --- an agent that receives a different answer on a retry will thrash --- and
it must answer fast enough that calling it is never the expensive option.
Clustering-based recovery satisfies neither constraint. Determinism here is a
design requirement, not an implementation detail.

\subsection{Delivering analysis an agent can afford}
Architecture analysis produces far more data than an agent can read. Serializing
the full dependency graph and recovered architecture of \arcade{} itself
--- 407 entities, 228 edges, 11 components --- costs \textbf{63{,}586
tokens}.\footnote{All token figures in this paper are estimates at
four characters per token, reproducible with \tool{measure\_budget.py} in the
paper's artifact directory.} Two mechanisms make that affordable.

\paragraph*{Session handles}
Every tool returns a compact summary carrying a \tool{session\_id}; the analysed
object itself stays server-side under a label (\tool{DependencyGraph},
\tool{Architecture}, \tool{SmellList}, \tool{MetricList}, \tool{IngestedRepo}).
Downstream tools accept either a session ID or an inline object, so an agent
pipes \tool{parse} into \tool{recover} into \tool{detect\_smells} while the graph
never crosses the wire. \tool{get\_full\_result} retrieves it only if the agent
genuinely needs it.

\paragraph*{Token budgeting}
Every tool additionally accepts \tool{max\_tokens}. Payloads containing a graph
or architecture are reduced by six ordered strategies, each re-checked against
the budget so the cheapest sufficient reduction wins: truncate component entity
lists; collapse entities to an FQN$\rightarrow$kind map; collapse edges to
per-relation counts; drop the entity map for kind counts; reduce components to
name and size; drop the package map. Everything else is handled generically by
repeatedly dropping the largest top-level key, preserving \tool{session\_id} and
flagging \tool{\_budget\_truncated} so the agent knows the view is partial.
Table~\ref{tab:budget} shows the effect on \arcade{}'s own graph.

\todo{Optional, and first on the cut list if space is tight, but it pre-empts a
reviewer who runs the numbers: the reduction is coarse. Between the 8k and 4k
budgets the payload falls from 6{,}468 to 216 tokens, because the fourth
strategy discards the entity map wholesale rather than sampling it. The budget
is respected but rarely approached from below.}

\subsection{Caching and enforcement}
Parsed graphs are cached under \tool{.arcade-cache/}. A content-hash-keyed
extract cache reuses the per-file pass for byte-identical files, so the
\tool{git checkout} and \tool{touch} churn of an agent session does not force a
re-parse; the link pass always runs in full, so edges can never go stale. This
incremental path is currently wired for the Python parser only.

\arcguard{} then enforces in two tiers. Advisory: the MCP tools, plus a
\tool{PostToolUse} hook that runs a check after each edit and injects findings
into the agent's context, so agents that forget to ask still get told. Blocking:
the same conformance engine as a pre-commit hook and a CI gate, exiting non-zero
on failure. Both tiers call one deterministic core, so the advisory answer an
agent receives mid-session and the verdict CI produces later cannot disagree.

\begin{figure}[t]
  \centering
  \fbox{\parbox[c][3.4cm][c]{0.92\columnwidth}{\centering
    \todo{FIGURE 2 --- system diagram: agent $\rightarrow$ MCP guardrail tools
    $\rightarrow$ deterministic conformance core, reading the spec and the
    baseline, over the \arcade{} engine; the verdict fans out to the advisory
    and blocking tiers. An ASCII draft exists under ``System shape'' in
    arcade-analyze-skill/GUARDRAIL\_PLAN.md --- redraw in TikZ.}}}
  \caption{\arcade{} in the agent's loop.}
  \label{fig:pipeline}
\end{figure}

% -----------------------------------------------------------------------------
% TABLE: measured, not estimated. Reproduce with measure_budget.py.
% -----------------------------------------------------------------------------
\begin{table}[t]
  \centering
  \caption{Token budgeting applied to \arcade{}'s own dependency graph and
           recovered architecture (63{,}586 tokens unreduced).}
  \label{tab:budget}
  \footnotesize
  \begin{tabular}{@{}rrl@{}}
    \toprule
    \textbf{Budget} & \textbf{Actual} & \textbf{Deepest reduction applied} \\
    \midrule
    --      & 63{,}586 & none \\
    16{,}000 & 12{,}644 & entities collapsed to FQN$\rightarrow$kind \\
     8{,}000 &  6{,}468 & edges collapsed to relation counts \\
     4{,}000 &     216 & entity and package maps dropped \\
    \bottomrule
  \end{tabular}
\end{table}

%==============================================================================
\section{The Tool in Action}
\label{sec:action}
%==============================================================================
The demonstration follows one agent session end to end. \todo{Keep this section
and the video in lockstep, beat for beat, so a reviewer can follow both at once.}

\paragraph*{Setup}
\tool{pip install arcade-agent[mcp,languages]} installs the library and the
\tool{arcade-mcp} console script, which the agent launches over stdio. No
project configuration is required to start reading a codebase; the guardrail
additionally needs an \tool{architecture.spec.json}, which \tool{init\_spec}
scaffolds from a template (layered, hexagonal, clean or MVC) or from a
one-paragraph description of intent.

\paragraph*{Understand}
The agent's first question about an unfamiliar repository is what is in it.
\tool{summarize} answers with the package tree, dependency hotspots, entry
points and entity-kind counts in a single call. On \arcade{} itself that
response is \textbf{911 tokens}. Reading the same source directly is roughly
\textbf{137{,}700 tokens} --- a hundred and fifty times more, and most of it
irrelevant to the question actually asked.

\paragraph*{Target}
Next the agent needs to know what to read. Given the task \emph{``add a Ruby
language parser''}, \tool{context\_for\_task} returns 14 ranked files in 3{,}615
tokens, each with a role and a reason. The four highest-scoring are the existing
Rust, Python, Kotlin and TypeScript parsers --- the files a contributor would in
fact open to write a new one --- recovered without the word \emph{parser} ever
being matched against a directory name. \todo{Ranking is lexical over entity
names, packages and paths, boosted by component membership. It degrades when the
task's vocabulary and the code's diverge; on a query phrased in terms the
codebase does not use, the right file is still in the ranked set but not at the
top. One sentence, honest, and worth more than the space it costs.}

\paragraph*{Guard}
Before writing, the agent calls \tool{propose\_placement} with a description of
what it is about to add and receives the component it belongs in together with
what that component may depend on. Before a cross-component import it calls
\tool{preview\_impact}, which reports the edges the change would introduce and
whether any violates the spec --- while the change is still hypothetical.
Fig.~\ref{fig:transcript} shows the moment that matters: the agent is told the
edge would violate, and revises its plan instead of its code.

\paragraph*{Gate}
After the edit, \tool{check\_architecture} returns a verdict, each violation
carrying the spec's own rationale and a concrete fix. The identical engine runs
as a pre-commit hook and in CI, where \tool{--fail-on error} exits non-zero and
the findings are posted to the pull request.

\paragraph*{Self-application}
\arcade{} analyses itself on every push. At v0.3.0 that run reports 76 source
files, 407 entities, 228 edges, 11 recovered components and 4 smells, and
commits \tool{.arcade/baseline.json} so the next run is compared against the
last good state rather than against nothing. The published report is the same
HTML artifact the \tool{visualize} tool produces for any project.

\begin{figure}[t]
\begin{lstlisting}
> propose_placement("Redis cache for cluster results")
  component: clustering
  may depend on: domain

> preview_impact(from="visualization", to="clustering")
  WOULD VIOLATE
  forbidden-dependency: visualization -> clustering
  why: "visualization must not drive analysis"
  fix: expose the result through the domain layer

> check_architecture()
  PASS  0 errors  0 new smells vs baseline
\end{lstlisting}
  \caption{An agent changing course before writing the code.}
  \label{fig:transcript}
\end{figure}

% -----------------------------------------------------------------------------
% All 19 names verified against origin/main @ 1b3532b (v0.3.0).
% NOTE: 19 counts MCP-registered tools (@server.tool()). Only 15 functions carry
% @tool -- ingest and parse are plain functions exposed over MCP. Never put "19"
% next to a number derived from list_tools(); say which surface you mean.
% -----------------------------------------------------------------------------
\begin{table}[t]
  \centering
  \caption{The 19 MCP tools, by purpose.}
  \label{tab:tools}
  \footnotesize
  \begin{tabular}{@{}p{0.24\columnwidth}p{0.68\columnwidth}@{}}
    \toprule
    \textbf{Group} & \textbf{Tools} \\
    \midrule
    Core pipeline &
      \tool{ingest}, \tool{parse}, \tool{analyze}, \tool{recover},
      \tool{detect\_smells}, \tool{compute\_metrics}, \tool{compare},
      \tool{query}, \tool{visualize} \\
    \addlinespace
    Comprehension &
      \tool{summarize}, \tool{explain\_component}, \tool{find\_relevant},
      \tool{api\_surface} \\
    \addlinespace
    Change-aware &
      \tool{diff\_impact}, \tool{dependency\_cone},
      \tool{changelog\_architecture}, \tool{context\_for\_task} \\
    \addlinespace
    Session &
      \tool{get\_full\_result}, \tool{list\_sessions} \\
    \bottomrule
  \end{tabular}
  \\[2pt]
  \todo{SPACE: Sections III--V now carry two tables and two figures. This is the
        first thing to cut (outline cut-list item 2) --- the measured budget
        table earns its space, an inventory does not. Collapse to a running
        two-column list in the text, or drop it.}
\end{table}

\todo{FIGURE 3 above is illustrative and MUST be replaced with a real captured
session before submission. The shape and the tool names are accurate, but the
transcript itself is reconstructed, and reviewers can tell. Capture it on
\tool{arcade\_core}, where the guardrail has already been verified to return
exactly this class of finding.}

%==============================================================================
\section{Evidence It Works}
\label{sec:evidence}
\secbudget{1.2 col + Table~\ref{tab:eval}}
%==============================================================================
\todo{Short, concrete, and scrupulous about scale. A demo paper does not need a
full study; it needs credible evidence the tool does what it claims. All numbers
below come from arcade-analyze-skill/evals/EVAL\_REPORT.md.}

\todo{(1) Detection is sound and deterministic: 3 of 3 injected violations
caught, each with the correct rule and a concrete fix; the gate exits non-zero
every time.}

\todo{(2) No false positives: across every run, each PASS was genuinely
conformant and each FAIL a genuine forbidden import, verified by inspecting the
imports.}

\todo{(3) Behavior changes exactly where theory predicts. REPORT THE NULL
RESULT -- on a pre-structured codebase there was no difference between
conditions, because agents imitate the layering already present. It is the
single most credibility-building sentence available, and it sharpens the claim:
the guardrail matters most on new code, where there is nothing to copy.}

\begin{table}[t]
  \centering
  \caption{Guardrail evaluation. Percentages are directional, not a benchmark.}
  \label{tab:eval}
  \footnotesize
  \begin{tabular}{@{}llrrr@{}}
    \toprule
    \textbf{Setting} & \textbf{Condition} & \textbf{Runs} & \textbf{Viol.} & \textbf{Rate} \\
    \midrule
    Pre-structured & without guardrail &  6 & 0 & 0\,\% \\
    Pre-structured & with guardrail    &  6 & 0 & 0\,\% \\
    \addlinespace
    Greenfield     & without guardrail & 14 & 3 & 21\,\% \\
    Greenfield     & with guardrail    & 14 & 0 & \phantom{0}0\,\% \\
    \addlinespace
    Injected       & detection         &  3 & 3 & 100\,\% caught \\
    \bottomrule
  \end{tabular}
  \\[2pt]
  \todo{Confirm the ``Viol.'' column header reads correctly for the injected
        row, where 3/3 means caught rather than committed -- consider splitting
        that row out or footnoting it.}
\end{table}

\subsection*{Threats to validity}
\todo{Three or four sentences, and do not skip them. All are already written
honestly in EVAL\_REPORT.md -- compress, do not soften. (a) Small $N$ and agent
stochasticity: the largest single cell is 10 runs per condition, so treat
percentages as directional. (b) Both conditions were told the folder roles, so
21\,\% is likely a floor for the unaided rate. (c) The more capable model
complied in every condition, so the measured benefit concentrates on
smaller and faster agents. (d) Conformance here is structural, not functional
correctness of the implemented features.}

%==============================================================================
\section{Related Work}
\label{sec:related}
\secbudget{0.6 col -- dense and citation-heavy}
%==============================================================================
\todo{Architecture recovery: ACDC, ARC, LIMBO, WCA, and ARCADE as a
workbench~\cite{tzerpos2000acdc,garcia2011arc,andritsos2004limbo,maqbool2004wca,laser2020arcade}.
Position \arcade{} as reimplementing and repackaging this line -- say it outright.}

\todo{Decay and conformance: decay studies and change-proneness
work~\cite{le2018decay,le2021decay-predictor,behnamghader2017evolution},
reflexion models~\cite{murphy1995reflexion}, and industrial conformance tools
such as ArchUnit~\cite{archunit}. Position: all post-hoc, human-invoked,
CI-time.}

\todo{Agent context and repository comprehension: repository maps, retrieval for
code agents, SWE-bench-style harnesses~\cite{TODO-agent-context}. Position:
these optimize \emph{what to read}; none carries an architectural contract.}

\todo{Guardrails in the agent loop: linters, type checkers, and tests already
sit there. Position: the missing rung is architecture. Cite the MCP
specification~\cite{mcp} for the integration surface.}

%==============================================================================
\section{Availability, Limitations, and Conclusion}
\label{sec:conclusion}
\secbudget{0.4 col}
%==============================================================================
\todo{Availability: \tool{pip install arcade-agent[mcp,languages]}; Python
$\ge$3.12; MIT. Repository \url{https://github.com/lemduc/arcade-agent};
documentation and usage instructions \url{https://arcade-agent.dev}; PyPI
\tool{arcade-agent} v0.3.0; video \url{TODO-video-url}. The track requires a
link to the publicly available tool AND its usage instructions -- both are
submission requirements, so verify both resolve from a clean machine.}

\todo{Limitations, one honest sentence each: cross-language relinking is
validated only on the JVM family; incremental parsing is wired for Python only,
so the sub-second guarantee does not yet generalize to very large polyglot
repos; MCP sessions are in-process and do not survive a server restart;
conformance is structural, not behavioral.}

\todo{Conclusion: one or two sentences. No new claims.}

\balance
\bibliographystyle{IEEEtran}
\bibliography{refs}

\end{document}
